% GENERATED FILE. Edit canonical lessons, then run npm run generate:books.
% canonical-source-hash: fnv1a64:d81e43aa242c2f74
% canonical-lessons: JA-C24-itami, JA-C24-ojiisan, JA-C24-ojisan, JA-C24-musuko, JA-C24-musume

\chapter{Pain, Grandfather, Uncle, Son, Daughter}
\label{ch:ja-pain-grandfather-uncle-son-daughter}
\hlchaptermodality{24}

\begin{chapteropening}
\textbf{By the end of this chapter:} \emph{I can say pain, a grandfather, an uncle, a son and a daughter.}

Complete the last lesson of chapter 24: I can say pain, a grandfather, an uncle, a son and a daughter.
\end{chapteropening}

\section[\texorpdfstring{\ja{いたみ} (itami)}{いたみ (itami)}]{\ja{いたみ} (itami) — pain}

\label{lesson:JA-C24-itami}



\begin{quote}
Before the new one: say the Japanese for fun, then the Japanese for scary.
\end{quote}

\subsection*{You'll want to know: \ja{いたみ}}
\textbf{\ja{いたみ}} — \emph{itami} — \textquotedblleft{}pain\textquotedblright{}.

\begin{cousinweb}
\textbf{\ja{いたみ}} (\emph{itami}) is native: the noun made from the verb \emph{itamu}, \textquotedblleft{}to hurt\textquotedblright{}, by its stem, beside the adjective \emph{itai}, \textquotedblleft{}painful\textquotedblright{} (JMdict). It is inherited rather than borrowed; its kanji is a kun'yomi attached later.
\end{cousinweb}

\begin{grammarlens}[title={What you've built}]
One more word for this chapter.
\end{grammarlens}

\subsection*{Your turn}
\begin{itemize}
\raggedright
  \item \emph{Say it:} \emph{itami}
  \item \emph{Say it:} \emph{itami}, once more
  \item \emph{Say it:} say \emph{itami}
  \item \emph{Recall:} say the Japanese for fun, then the Japanese for scary, then say \emph{itami} again
\end{itemize}

\begin{tcolorbox}[breakable,colback=teal!4,colframe=teal!35!black,title={Before you move on}]
What does \ja{いたみ} mean? (\textquotedblleft{}Pain\textquotedblright{}.) Say it once more.
\end{tcolorbox}
\section[\texorpdfstring{\ja{おじいさん} (ojiisan)}{おじいさん (ojiisan)}]{\ja{おじいさん} (ojiisan) — a grandfather}

\label{lesson:JA-C24-ojiisan}



\begin{quote}
Before the new one: say the Japanese for scary, then the Japanese for pain.
\end{quote}

\subsection*{You'll want to know: \ja{おじいさん}}
\textbf{\ja{おじいさん}} — \emph{ojiisan} — \textquotedblleft{}a grandfather\textquotedblright{}.

\begin{cousinweb}
\textbf{\ja{おじいさん}} (\emph{ojīsan}) is native \emph{jii}, \textquotedblleft{}old man; grandfather\textquotedblright{}, with the polite prefix \emph{o-} and the polite ending \emph{-san}. JMdict writes it with characters for \textquotedblleft{}grandfather\textquotedblright{}, read for their meaning. It is inherited rather than borrowed.
\end{cousinweb}

\begin{grammarlens}[title={What you've built}]
One more word for this chapter.
\end{grammarlens}

\subsection*{Your turn}
\begin{itemize}
\raggedright
  \item \emph{Say it:} \emph{ojiisan}
  \item \emph{Say it:} \emph{ojiisan}, once more
  \item \emph{Say it:} say \emph{ojiisan}
  \item \emph{Recall:} say the Japanese for scary, then the Japanese for pain, then say \emph{ojiisan} again
\end{itemize}

\begin{tcolorbox}[breakable,colback=teal!4,colframe=teal!35!black,title={Before you move on}]
What does \ja{おじいさん} mean? (\textquotedblleft{}A grandfather\textquotedblright{}.) Say it once more.
\end{tcolorbox}
\section[\texorpdfstring{\ja{おじさん} (ojisan)}{おじさん (ojisan)}]{\ja{おじさん} (ojisan) — an uncle}

\label{lesson:JA-C24-ojisan}



\begin{quote}
Before the new one: say the Japanese for pain, then the Japanese for a grandfather.
\end{quote}

\subsection*{You'll want to know: \ja{おじさん}}
\textbf{\ja{おじさん}} — \emph{ojisan} — \textquotedblleft{}an uncle\textquotedblright{}.

\begin{cousinweb}
\textbf{\ja{おじさん}} (\emph{ojisan}) is native \emph{oji}, \textquotedblleft{}uncle\textquotedblright{}, with the polite prefix \emph{o-} and the ending \emph{-san}. JMdict writes \emph{oji} with two characters meaning \textquotedblleft{}father's elder (or younger) brother\textquotedblright{}, read as a whole for their meaning. It is inherited rather than borrowed.
\end{cousinweb}

\begin{grammarlens}[title={What you've built}]
One more word for this chapter.
\end{grammarlens}

\subsection*{Your turn}
\begin{itemize}
\raggedright
  \item \emph{Say it:} \emph{ojisan}
  \item \emph{Say it:} \emph{ojisan}, once more
  \item \emph{Say it:} say \emph{ojisan}
  \item \emph{Recall:} say the Japanese for pain, then the Japanese for a grandfather, then say \emph{ojisan} again
\end{itemize}

\begin{tcolorbox}[breakable,colback=teal!4,colframe=teal!35!black,title={Before you move on}]
What does \ja{おじさん} mean? (\textquotedblleft{}An uncle\textquotedblright{}.) Say it once more.
\end{tcolorbox}
\section[\texorpdfstring{\ja{むすこ} (musuko)}{むすこ (musuko)}]{\ja{むすこ} (musuko) — a son}

\label{lesson:JA-C24-musuko}



\begin{quote}
Before the new one: say the Japanese for a grandfather, then the Japanese for an uncle.
\end{quote}

\subsection*{You'll want to know: \ja{むすこ}}
\textbf{\ja{むすこ}} — \emph{musuko} — \textquotedblleft{}a son\textquotedblright{}.

\begin{cousinweb}
\textbf{\ja{むすこ}} (\emph{musuko}) is a native word, inherited rather than borrowed. It is usually explained as \emph{musu}, \textquotedblleft{}to be born, to grow\textquotedblright{}, with \emph{ko}, \textquotedblleft{}child\textquotedblright{}; JMdict's kanji, \textquotedblleft{}breath\textquotedblright{} and \textquotedblleft{}child\textquotedblright{}, are read as a whole for their meaning.
\end{cousinweb}

\begin{grammarlens}[title={What you've built}]
One more word for this chapter.
\end{grammarlens}

\subsection*{Your turn}
\begin{itemize}
\raggedright
  \item \emph{Say it:} \emph{musuko}
  \item \emph{Say it:} \emph{musuko}, once more
  \item \emph{Say it:} say \emph{musuko}
  \item \emph{Recall:} say the Japanese for a grandfather, then the Japanese for an uncle, then say \emph{musuko} again
\end{itemize}

\begin{tcolorbox}[breakable,colback=teal!4,colframe=teal!35!black,title={Before you move on}]
What does \ja{むすこ} mean? (\textquotedblleft{}A son\textquotedblright{}.) Say it once more.
\end{tcolorbox}
\section[\texorpdfstring{\ja{むすめ} (musume)}{むすめ (musume)}]{\ja{むすめ} (musume) — a daughter}

\label{lesson:JA-C24-musume}



\begin{quote}
Before the new one: say the Japanese for an uncle, then the Japanese for a son.
\end{quote}

\subsection*{You'll want to know: \ja{むすめ}}
\textbf{\ja{むすめ}} — \emph{musume} — \textquotedblleft{}a daughter\textquotedblright{}.

\begin{cousinweb}
\textbf{\ja{むすめ}} (\emph{musume}) is a native Japanese word, inherited rather than borrowed from Chinese, and its older origin is not settled. Its kanji was attached to it later as a kun'yomi, a native reading given to a Chinese character (JMdict, KANJIDIC2).
\end{cousinweb}

\begin{grammarlens}[title={What you've built}]
One more word for this chapter.
\end{grammarlens}

\subsection*{Your turn}
\begin{itemize}
\raggedright
  \item \emph{Say it:} \emph{musume}
  \item \emph{Say it:} \emph{musume}, once more
  \item \emph{Say it:} say \emph{musume}
  \item \emph{Recall:} say the Japanese for an uncle, then the Japanese for a son, then say \emph{musume} again
\end{itemize}

\begin{tcolorbox}[breakable,colback=teal!4,colframe=teal!35!black,title={Before you move on}]
What does \ja{むすめ} mean? (\textquotedblleft{}A daughter\textquotedblright{}.) Say it once more.
\end{tcolorbox}
